\minitoc

% ========================================================================================================
% Introduction 

\section{Introduction}

\subsection{Motivations} 

En 2011, Daniel Kahneman publie un livre intitulé \emph{Système 1 / Système 2 : Les deux vitesses de la pensée} 
dans lequel il présente deux modes de pensée : les systèmes 1 et 2 (S1 et S2). 
Le Système 1 est comparable à l'intuition. Il ne repose pas sur des bases théoriques complexes et 
est difficilement explicable. Il résulte d'un long apprentissage et d'expérimentations. 
Le Système 2, quant à lui, est constitué d'un raisonnement construit pas à pas et suivant 
une logique connue. Il est, de fait, plus lent que le Système 1, mais explicable. 

L'Intelligence Artificielle actuelle, notamment les réseaux de neuronnes, peuvent être vus comme des 
Systèmes 1. En effet, ils sont entraînés sur beaucoup de données étiquetées (les expériences)
et leur fonctionnement est difficilement explicable lors de la phase d'inférence. 

L'objectif de ce cours et de développer une logique propositionnelle permettant de formaliser 
de tels raisonnements et les rendre explicables. Plus globablement, l'enjeu de ce domaine de 
recherche est de développer une IA hybride mélangeant système 1 et 2. 

\subsection{Différentes imperfections de l'information}

Les informations disponibles sur un système (météo, population animale, etc...) sont généralement 
imparfaites. Nous devons donc être capables de définir cette imperfection. L'objectif de ce cours est 
donc de la quantifier et d'en présenter différents formalismes de représentation. 
On distingue ainsi 4 natures d'imperfection de l'information : 

\begin{itemize}
    \item \textbf{L'incohérence} : des connaissances contradictoires. 
    \item \textbf{L'imprécision} : des connaissances mal définies pour lesquelles on donnera un degré 
    de vérité. Cela peut aussi être des disjonctions ou des connaissances incomplètes.
    \item \textbf{L'incertitude} : des connaissances dont on a un doute sur la source. 
    \item \textbf{La non-pertinence} : les informations données ne sont pas en lien avec le sujet traité. 
\end{itemize}

\begin{example}
    Soit la question et les réponses suivantes : 
    \begin{quote}
        \emph{Où avez-vous perdu votre téléphone ? }

        \begin{enumerate}
            \item "Près de l'amphi Maxwell." (imprécis/vague)
            \item "Je suis presque certain que c'est en B08." (incertain)
            \item "Soit dans Maxwell, soit dans B08." (imprécis)
            \item "Je crois que c'est près de Maxwell." (incertain/vague)
        \end{enumerate}
    \end{quote}
\end{example}

Il y a une relation entre incertain et imprécis : moins on est précis, plus on est certain. 

\begin{remark}
    La théorie des probabilités permet de quantifier mathématiquement cette incertitude. 
    Une bonne connaissance de celle-ci s'apparente à connaître la loi de tous les évènnements possibles.
    En revanche, pour la connaître, il est nécessaire de pouvoir répéter un certain nombre de fois 
    l'expérience aléatoire, ce qui n'est pas toujours le cas. 
\end{remark}

% ========================================================================================================
% Mesures d'incertitude

\section{Mesures d'incertitude}

Nous allons ici étendre la théorie des probabilités en définissant des mesures d'incertitudes. 
On suppose qu'il n'est, a priori, pas possible de répéter plusieurs fois l'expérience aléatoire. 
Nous devons donc nous contenter de \emph{connaissances incertaines}.

\subsection{Élicitation des probabilités par les paris}

Soit $ \Omega$ un univers et $ u : \Omega \longrightarrow \R$ définissant l'utilité d'être dans un 
état pour un agent prenant les décisions $f$. On définit $ \forall s, s' \in \Omega, s' = f(s)$ le changement 
d'état entre $s'$ et $s$. On suppose qu'il existe une distribution de probabilités 
$ \myP : \Omega \longrightarrow [0,1]$. 

\begin{definition}[Utilité Espérée]
    L'intérêt d'une décision $f$ est définie par la somme des utilités de ses conséquences
    pondérées par la probabilité que chaque état se réalise : 
        \begin{equation}
            UE(f) = \sum_{ \omega \in \Omega} \myP(s) u(f(s)) 
        \end{equation}
\end{definition}

Or ici, on suppose l'existence d'une probabilité. Définissons la notion de papi pour nous 
permettre de la retrouver. 

\begin{definition}[Pari]
    Soit $A \subset \Omega$ un évènnement et $y,x \in \R$. On définit le pari $yAx$ comme 
    "on gagne $y$ si $A$ et $x$ sinon". 

    L'utilité espérée d'un pari $yAx$ est définie par : 
        $$ UE(yAx) = u(y) \myP(A) + u(x) \myP( \overline{A}). $$ 
\end{definition}

Soient $A,B \subset \Omega$ et $ y \in \R$. Dire qu'un agent croit plus en $A$ qu'en $B$ 
revient à dire qu'il préfère parier $yA0$ plutôt que $yB0$. Autrement dit : 
    \begin{align*}
       & UE(yA0) > UE(yB0) \\ 
       \iff & u(y) \myP(A) u(0) > u(y) \myP(B) u(0) \\ 
       \iff & \myP(A) > \myP(B) 
    \end{align*}

On cherche maintenant à estimer la valeur numérique de $ \myP(A)$. Pour cela, on demande à un agent 
"Quelle somme $x$ accepterais-tu de payer pour renoncer au pari $yA0$ ?"
Lorsque l'agent hésite, on a : 
    $$ u(x) = UE(yA0) = \myP(A) u(y) $$ 
    $$ \iff \myP(A) = \frac{u(x)}{u(y)} $$ 
Ainsi la probabilité d'un évènnement $A$ peut être estimée par le rapport des utilités entre ce que l'agent 
est prêt à engager $x$ et ce qu'il peut gagner $y$. 
La \emph{probabilité subjective} est donc la somme relative que l'on est prêt à parier. 

\subsection{Probabilités imprécises}

Ce nouveau modèle n'est en revanche pas suffisant, l'exemple suivant nous le montre. 

\begin{proposition}[Paradoxe d'Ellsberg]
    Soit une urne contenant 90 boules dont 30 rouges et 60 noires ou jaunes. 
    On ne sait pas quelle est la proportion de boules noire parmi les noires et jaunes. 
    On effectue le tirage d'une boule et on définit différents paris : 
    \begin{itemize}
        \item Option A : recevoir 100 euros si la boule est rouge, rien sinon.
        \item Option B : recevoir 100 euros si la boule est noire, rien sinon. 
        \item Option C : recevoir 100 euros si la boule est rouge ou jaune, rien sinon. 
        \item Option D : recevoir 100 euros si la boule noire ou jaune, rien sinon. 
    \end{itemize}
    On constate que les agents préfèrent parier sur Rouge ($ \myP = 1/3$) plutôt que sur Noir (de probabilité inconnue). 
    Dans le même temps, ils préfèrent parier sur "Noir ou Jaune" plutôt que sur "Rouge ou Jaune", ce qui constitue 
    une incohérence. 
    Il n'existe donc aucune distribution de probabilité unique capable d'expliquer ces choix simultanés via l'UE.
\end{proposition}

Ainsi le modèle de l'utilité espérée ne suffit pas toujours à décrire le comportement humain face à une ignorance 
totale. 
Lorsque l'information est incomplète (comme dans le paradoxe d'Ellsberg), il est impossible de définir une probabilité 
unique $p$. On travaille alors avec un ensemble $\mathcal{F}$ de mesures de probabilité compatibles avec 
l'information disponible.

\begin{definition}[Famille de probabilités]
    On définit $ \mathcal{F}$ comme l'ensemble des mesures de probabilités $p$ compatibles 
    avec l'expérience aléatoire (i.e avec les données de l'énoncé). 
\end{definition}

\begin{example}[Urnes d'Ellsberg]
    Dans notre exemple précédent, cet ensemble est défini par : 
        $$ \mathcal{F} := \{p : \Omega \longrightarrow [0,1] \mid p(R) = 1/2, p(N) + p(J) = 2/3\} $$
\end{example}

\begin{definition}[Mesures infra et supra-probabilistes]
    Soit $\mathcal{F}$ un ensemble de distributions de probabilité. 
    Pour tout événement $A \subseteq \Omega$, on définit :
    \begin{itemize}
        \item \textbf{La probabilité inférieure :} $P_*(A) = \inf_{p \in \mathcal{F}} P(A)$.
        \item \textbf{La probabilité supérieure :} $P^*(A) = \sup_{p \in \mathcal{F}} P(A)$.
    \end{itemize}
\end{definition}

Face à l'imprécision, le critère de décision consiste à maximiser l'\emph{utilité espérée minimale} (approche pessimiste) 
définie par : 
\begin{equation}
    UE_*(f) = \inf_{p \in \mathcal{F}} UE(f)
\end{equation}

\begin{prop}[Mesures infra/supra-probabilistes]
    Ces probabilités vérifient :
    \begin{itemize}
        \item \emph{Encadrement :} $\forall A \subset \Omega, \quad P_*(A) \leqslant P(A) \leqslant P^*(A)$.
        \item \emph{Dualité :} $P_*(A) = 1 - P^*(\overline{A})$.
        \item \emph{Super-additivité :} $P_*$ est super-additive : $P_*(A \cup B) \geqslant P_*(A) + P_*(B)$ (pour $A, B$ disjoints).
        \item \emph{Sub-additivité :} $P^*$ est sub-additive : $P^*(A \cup B) \leqslant P^*(A) + P^*(B)$ (pour $A, B$ disjoints).
    \end{itemize}
\end{prop}

\begin{remark}
    Cette approche permet de différencier la notion d'ignorance de l'équiprobabilité. 
    En probabilités classiques, lorsque l'on ne connaît pas les données de l'expérience aléatoire, on suppose que 
    l'on se trouve en situation d'équiprobabilité. L'exemple des urnes nous en montre la faiblesse. 
    Les mesures définies ci-dessus permettent donc de palier à ce problème et de différencier des deux notions. 
\end{remark}

\begin{example}[Urnes d'Ellsberg]
    Soit $ \mathcal{F}$ une famille de probabilités pour ce problème. On choisira maintenant la décision 
    $f$ qui maximise $UE_*(f) = \inf_{p \in \mathcal{F}} UE(f)$. On aura donc : 
        \begin{align*}
            UE_*(100\{\text{Rouge}\}0) &= u(100) \times \myP(\{ \text{Rouge} \}) = \frac{1}{3} u(100) \\ 
            UE_*(100 \{ \text{Noir} \}0) &= \min_{p \in \mathcal{F}} p(\{ \text{Noir} \}) \times u(100) = 0 \\
            UE_*(100 \{ \text{Rouge, Jaune} \} 0) &= \min_{p \in \mathcal{F}} p(\{ \text{Rouge, Jaune} \}) u(100) = \frac{1}{3} u(100) \\ 
            UE_*(100 \{ \text{Jaune, Noir} \} 0) &=  \min_{p \in \mathcal{F}} p(\{ \text{Noir, Jaune} \}) u(100) = \frac{2}{3} u(100)
        \end{align*}
    On a donc : 
        $$ UE_*(100\{\text{Rouge}\}0) > UE_*(100 \{ \text{Noir} \}0) $$ 
        $$ UE_*(100 \{ \text{Rouge, Jaune} \} 0) > UE_*(100 \{ \text{Jaune, Noir} \} 0). $$ 
    Ce qui concorde avec les choix "humains". 
\end{example}

\subsection{Fonction de croyance - Théorie de Dempster-Shafer}

Pour respecter ce nouveau cadre, nous devons donc définir une nouvelle structure probabiliste. 
Celle-ci s'appuiera sur les \emph{fonctions de masse}. 

\begin{definition}[Fonction de masse]
    Soit $ \Omega$ un univers. On définit une fonction de masse sur $ \Omega$ comme une application : 
        \begin{equation}
            m = \mathcal{P}( \Omega) \longrightarrow [0,1] 
        \end{equation}
    qui vérifie : 
    \begin{enumerate}[label=\roman*)]
        \item $ m(\emptyset) = 0 $ 
        \item $ \displaystyle \sum_{A \in \mathcal{P}( \Omega)} m(A) = 1$
    \end{enumerate}
    Cette application définit la probabilité que l'ensemble $A$ soit l'état actuel. Lors que $m(A) > 0$, on dit 
    que $A$ est un \emph{élément focal}. 
\end{definition}

\begin{definition}[Croyance et Plausibilité]
    Soit $ \Omega$ un univers. On définit deux applications : 
    \begin{itemize}
        \item \textbf{La croyance} $Bel$ définie par : 
            \begin{equation}
                Bel : 
                    \begin{cases}
                        \mathcal{P}( \Omega) \longrightarrow \R \\ 
                        A \longmapsto \displaystyle \sum_{B \in \mathcal{P}( \Omega) \mid B \subset A} m(B) 
                    \end{cases}
            \end{equation}
            qui estime dans quelle mesure un évènement $A$ est certain. 
        \item \textbf{La plausibilité} $pl$ définie par : 
            \begin{equation}
                pl : 
                    \begin{cases}
                        \mathcal{P}( \Omega) \longrightarrow \R \\ 
                        A \longmapsto \displaystyle \sum_{B \in \mathcal{P}( \Omega) \mid B \cap A \not = 0} m(B).
                    \end{cases}
            \end{equation}
            qui détermine en quelle mesure $A$ est plausible. 
    \end{itemize}
\end{definition}

\begin{prop}[Croyance et Plausibilité]
    Ces deux applications vérifient : 
    \begin{equation}
        \forall A \in \mathcal{P}( \Omega), \quad pl(A) = 1 - Bel( \overline{A}) 
    \end{equation}
    et : 
    \begin{equation}
        \forall A \in \mathcal{P}( \Omega), \quad Bel(A) \leqslant pl(A)
    \end{equation}
    Un évènnement est donc toujours plus plausible que certain. 
\end{prop}

Ces deux nouvelles notions nous permettent de détailler la différence entre équiprobabilité et 
ignorance. 

\begin{example}
    Soient trois chevaux $c_1, c_2$ et $c_3$. On veut savoir lequel des trois gagnera une course. 
    Deux experts donnent leur avis sur son issue : 
    \begin{itemize} 
        \item Expert 1 : "Chacun des trois chevaux a la même chance de gagner."
        \item Expert 2 : "Je n'en sais rien." 
    \end{itemize}
    En théorie des probabilités classiques, on aurait donc : 
    \begin{itemize} 
        \item Expert 1 : $ \myP(c_1) = \myP(c_2) = \myP(c_3) = 1/3$ 
        \item Expert 2 : $ \myP(c_1) = \myP(c_2) = \myP(c_3) = 1/3$ 
    \end{itemize}
    Mais en théorie de fonction de croyance : 
    \begin{itemize} 
        \item Expert 1 : $ m_1(c_1) = m_1(c_2) = m_1(c_3) = 1/3$ et 
            $$ 
                \begin{cases}
                    pl(\{"c_1 \text{ gagne}\}) = Bel(\{"c_1 \text{ gagne}\}) = 1/3 \\ 
                    pl(\{"c_2 \text{ gagne}\}) = Bel(\{"c_2 \text{ gagne}\}) = 1/3 \\ 
                    pl(\{"c_3 \text{ gagne}\}) = Bel(\{"c_3 \text{ gagne}\}) = 1/3 \\ 
                \end{cases}
            $$ 
        \item Expert 2 : $ m_2(\{c_1, c_2, c_3\}) = 1$ et 
            $$ 
                \begin{cases}
                    pl(\{"c_1 \text{ gagne}\}) = 1, Bel(\{"c_1 \text{ gagne}\}) = 0 \\ 
                    pl(\{"c_2 \text{ gagne}\}) = 1, Bel(\{"c_2 \text{ gagne}\}) = 0 \\ 
                    pl(\{"c_3 \text{ gagne}\}) = 1, Bel(\{"c_3 \text{ gagne}\}) = 0 \\ 
                \end{cases}
            $$
    \end{itemize}
\end{example}

\begin{remark}
    Les mesures de probabilités sont un cas particulier des mesures de croyance. En effet, si tous les éléments 
    focaux de $ \mathcal{P}( \Omega)$ sont des singletons, on retrouve une telle mesure. 
\end{remark}

\begin{proposition}
    En réalité, les fonctions de croyance $pl$ et $Bel$ sont des manières compactes de représenter 
    une famille de probabilités. On a ainsi : 
    \begin{itemize}
        \item $Bel$ correspond exactement à la probabilité inférieure $ \myP_*$ sur $ \mathcal{F}$. 
        \item $pl$ correspond aussi à la probabilité inférieure $ \myP^*$ sur $ \mathcal{F}$. 
    \end{itemize}
    La famille $ \mathcal{F}$ peut donc être réécrite : 
        $$ \mathcal{F} = \{p \mid \forall A \in \mathcal{P}( \Omega), \quad Bel(A) \leqslant p(A) \leqslant Pl(A)\} $$
\end{proposition}
 
On peut donc définir plus formellement l'utilité pessimiste. 

\begin{definition}[Utilité pessimiste]
    L'utilité pessimiste d'une décision $f$ est définie par :
    \begin{equation}
        \boxed{ U_{pes}(f) = \sum_{A \subseteq \Omega} m(A) \cdot \min_{s \in A} u(f(s)) }
    \end{equation}
\end{definition}

\begin{remark}
    Cette approche permet de retrouver les résultats du paradoxe d'Ellsberg de manière directe en utilisant les masses $m(\{Rouge\}) = 1/3$ et $m(\{Noir, Jaune\}) = 2/3$.
\end{remark}

\subsection{Théorie des possibilités}

La théorie des possibilités est particulièrement adaptée lorsque les connaissances sont représentées par des ensembles 
emboîtés ($E_1 \subseteq E_2 \subseteq \dots \subseteq E_n$). Elle repose sur une distribution de
possibilité $\pi$.

\begin{definition}[Distribution de possibilité]
    Soit $\pi : \Omega \longrightarrow [0,1]$ une application où $\pi(s)$ estime dans quelle mesure il est 
    possible que $s$ soit l'état réel.
    \begin{itemize}
        \item $\pi(s) = 1$ signifie que l'état $s$ est tout à fait possible (non surprenant).
        \item $\pi(s) = 0$ signifie que l'état $s$ est jugé impossible.
        \item $E_1 = \{s \mid \pi(s) = 1\}$ est le \textbf{noyau} de $\pi$.
        \item $E_n = \{s \mid \pi(s) > 0\}$ est le \textbf{support} de $\pi$.
    \end{itemize}
\end{definition}

Pour évaluer la confiance en un événement $A \subseteq \Omega$, on définit deux mesures duales :

\begin{definition}[Mesures de Possibilité et de Nécessité]
    \begin{itemize}
        \item \textbf{La mesure de possibilité $\Pi(A)$} s'appuie sur le cas le plus favorable à $A$ :
            \begin{equation}
                \boxed{\Pi(A) = \sup_{s \in A} \pi(s)}
            \end{equation}
        \item \textbf{La mesure de nécessité $N(A)$} s'appuie sur le cas le plus favorable de l'événement contraire :
            \begin{equation}
                \boxed{N(A) = 1 - \Pi(\overline{A}) = 1 - \sup_{s \notin A} \pi(s)}
            \end{equation}
    \end{itemize}
\end{definition}

\begin{prop}[Axiomes caractéristiques]
    Contrairement aux probabilités, ces mesures ne sont pas additives, mais respectent les règles suivantes :
    \begin{equation}
        \forall A, B \in \mathcal{P}( \Omega), \quad 
        \begin{cases}
            \Pi(A \cup B) = \max(\Pi(A), \Pi(B)) \\ 
            N(A \cap B) = \min(N(A), N(B))
        \end{cases}
    \end{equation}
\end{prop}

\begin{remark}[Lien avec Dempster-Shafer]
    La théorie des possibilités est un cas particulier des fonctions de croyance où les éléments focaux sont 
    \textbf{consonants} (emboîtés). Dans ce cas, la croyance $Bel$ coïncide avec la nécessité $N$, 
    et la plausibilité $pl$ avec la possibilité $\Pi$.
\end{remark}

\subsection{Conclusion : Panorama sur les mesures d'incertitude}

Toutes les mesures étudiées ($P, Bel, pl, N, \Pi$) s'inscrivent dans le cadre plus général 
des \emph{capacités}. 

\begin{definition}[Capacité]
    Une application $g : \mathcal{P}( \Omega) \longrightarrow [0,1]$ est une capacité si elle vérifie :
    \begin{itemize}
        \item \emph{L'impossibilité} : $g(\emptyset) = 0$. 
        \item \emph{La Masse Unitaire} : $g(\Omega) = 1$. 
        \item \emph{La Monotonie} : $A \subseteq B \Rightarrow g(A) \leqslant g(B)$. 
    \end{itemize}
\end{definition}

\begin{table}[h]
    \centering
    \renewcommand{\arraystretch}{1.5} % Améliore l'espacement vertical
    \begin{tabular}{@{}>{\bfseries}l c l@{}}
        \toprule
        \text{Mesure} & \text{Symbole} & \text{Axiome / Propriété caractéristique} \\
        \midrule
        Probabilité & $ \myP$ & $ \myP(A \cup B) = \myP(A) + \myP(B)$ (si $A \cap B = \emptyset$) \\
        Croyance & $Bel$ & $Bel(A \cup B) \geqslant Bel(A) + Bel(B) - Bel(A \cap B)$ \\
        Plausibilité & $Pl$ & $Pl(A \cup B) \leqslant Pl(A) + Pl(B) - Pl(A \cap B)$ \\
        Possibilité & $\Pi$ & $\Pi(A \cup B) = \max(\Pi(A), \Pi(B))$ \\
        Nécessité & $N$ & $N(A \cap B) = \min(N(A), N(B))$ \\
        \bottomrule
    \end{tabular}
    \caption{Panorama des mesures d'incertitude et leurs axiomes caractéristiques}
    \label{tab:synthese_mesures}
\end{table}

% ========================================================================================================
% Logique possibiliste 

\section{Logique possibiliste}

La logique classique est binaire : une proposition est soit vraie, soit fausse. Or, les connaissances d'un agent 
sont souvent entachées d'incertitude. La \emph{logique possibiliste} étend la logique classique en associant à 
chaque formule un degré de certitude (nécessité).

\subsection{Syntaxe et Sémantique}

\begin{definition}[Base de connaissances possibiliste]
    Une base de connaissances possibiliste $\Sigma$ est un ensemble de \emph{formules de logique 
    possibiliste} $(\phi_i, \alpha_i)$ telles que :
    \begin{itemize}
        \item $\phi_i$ est une formule logique propositionnelle classique.
        \item $\alpha_i \in ]0, 1]$ est le \textbf{degré de nécessité} (ou de certitude) associé à $\phi_i$.
    \end{itemize}
    On note $\Sigma = \{ (\phi_i, \alpha_i), i=1, \dots, n \}$.
\end{definition}

$ \alpha_i $ peut être interprété par une borne inférieure du degré de certitude d'une proposition 
$ \varphi_i$ telle que $ N( \varphi) > \alpha$, avec $N$ une mesure de nécessité. 

\begin{example}[Base de connaissances type]
    Soit la base $\Sigma$ concernant un oiseau $O$, un manchot $M$ et la capacité de voler $V$ définie par : 
    \begin{enumerate}
        \item $(M \to O, 1)$ : Il est certain qu'un manchot est un oiseau.
        \item $(O \to V, 0.8)$ : Il est assez certain qu'un oiseau vole.
        \item $(M \to \neg V, 0.9)$ : Il est très certain qu'un manchot ne vole pas.
    \end{enumerate}
\end{example}

Une base possibiliste $ \Sigma$ induit une base de logique classique, notée $ \Sigma^*$, définie par : 
    \begin{equation}
        \Sigma^* := \{ \varphi \mid \exists \alpha \in [0,1], ( \varphi, \alpha) \in \Sigma\}
    \end{equation}
Ainsi $ \varphi \in \Sigma^*$ définit un ensemble de mondes la vérifiant que l'on 
note $[ \varphi]$ définit par : 
    \begin{equation}
        [ \varphi] := \{ \omega \in \Omega \mid \omega \models \varphi\}
    \end{equation}

\begin{definition}[Distribution de possibilité induite]
    Une base de possibilité $ \Sigma$ induit une distribution de possibilité $ \pi$ 
    sur l'ensemble des mondes possibles $ \Omega$. On a donc : 
    \begin{equation}
        \pi_\Sigma(\omega) = 1 - \max \{ \alpha_i \mid \omega \not\models \phi_i \}
    \end{equation}
    Pour un monde $ \omega \in \Omega$, on a donc 3 cas différents : 
    \begin{itemize}
        \item si $ \omega$ ne contredit aucunes formules possibilistes de $ \Sigma$, alors 
        $ \pi ( \omega) = 1$. 
        \item si $ \omega$ contredit une connaissance certaine, alors $ \pi ( \omega) = 0$. 
        \item sinon $ \pi( \omega) \in [0, 1]$. 
    \end{itemize}
    \begin{quote}
        \emph{"Plus un monde viole une formule certaine, moins il est possible."}
    \end{quote}
\end{definition}

\begin{example}[Calcul de distribution de possibilité]
    Considérons le monde $\omega = (M, O, V)$ (un manchot oiseau qui vole). 
    Ce monde viole la formule $\phi_3 : M \to \neg V$ de degré $\alpha_3 = 0.9$. 
    En appliquant la sémantique de la base :
    \[ \pi_\Sigma(\omega) = 1 - \max \{ 0.9 \} = 0.1 \]
    Ce monde est donc jugé peu possible, car il contredit une connaissance quasi-certaine.
\end{example}

\subsection{Raisonnement}

On souhaite maintenant raisonner sur des bases possibilistes $ \Sigma$ dans le but d'en déduire des mondes 
possibles. 
Une première idée serait de calculer $ N( \omega)$ et $ \Pi( \omega) $ pour tout monde $ \omega \in \Omega$. 
Or, c'est un problème très coûteux en temps et en mémoire. 
Une seconde idée est d'essayer de se ramener au cas de la logique classique à partir de certaines réductions. 
C'est ce que nous allons détailler ici. 

\begin{notation}
    Pour de tels raisonnements, on introduit la notation $ \Sigma_{ \alpha} \models \varphi$. 
    Elle signifie que la base $ \Sigma$ permet de déuire que $ N ( \varphi) \geqslant \alpha$. 
\end{notation}

L'outil fondamental pour nous ramener au cas "simple" de la logique classique est la \emph{coupure de 
niveau}. 

\begin{definition}[Coupure de niveau]
    Soit $ \Sigma$ une base possibiliste. On définit la \emph{coupure de niveau de degré $ \alpha$}
    de $ \Sigma$ comme l'ensemble : 
        \begin{equation}
            \Sigma_\alpha := \{ ( \varphi, \alpha) \in \Sigma \mid \alpha_i \geqslant \alpha\}
        \end{equation}
\end{definition}

\begin{theorem}[Fondamental]
    Soit $ \Sigma$ une base possibiliste et $ ( \varphi, \alpha) \in \Sigma$. On a : 
        \begin{equation}
            N( \varphi) \geqslant \alpha \quad \text{ssi} \quad ( \Sigma_\alpha)^* \vdash \varphi 
        \end{equation}
\end{theorem}

Autrement dit, pour affirmer que $ N( \varphi) \geqslant \alpha$, il faut prouver 
mathématiquement que $ \varphi$ est vraie seulement à partir des formules de $ (\Sigma_\alpha)^*$, 
en ignorant les degrés. 

Pour cela, on définit un ensemble d'axiomes, desquels découleront des formules plus complexes. 

\begin{definition}[Axiomes de la logique possibiliste]
    Soient $ p,q$ deux formules propositionnelles. On a les axiomes suivants : 
    \begin{enumerate}[label=\roman*)]
        \item \emph{Masse unitaire} : $N(\top) = 1$ 
        \item \emph{Nécessaité} : $N( p \wedge q) = \min( N(p), N(q)) $
        \item \emph{Monotonie} : Si $ p \models q$, alors $ N(p) \leqslant N(q)$. 
    \end{enumerate}
    De même pour la forme duale : 
    \begin{enumerate}[label=\roman*)]
        \item \emph{Impossibilité} : $ \Pi( \bot ) = 0$ 
        \item \emph{Possibilité} : $ \Pi( p \vee q) = \max( \Pi(p), \Pi(q))$ 
        \item \emph{Monotonie} : Si $ p \models q$, alors $ \Pi(p) \leqslant \Pi(q)$. 
    \end{enumerate}
\end{definition}

\begin{proposition}
    Il existe différentes règles d'inférence. En voici quelques-unes :
    \begin{itemize}
        \item \emph{Modus Ponens possibiliste} : 
        \begin{equation}
            \text{Si } N( p \to q) \geqslant \alpha \text{ et } N(p) \geqslant \beta \text{ alors } N(q) \geqslant \min( \alpha, \beta)
        \end{equation}
        \begin{equation}
            \text{Si } N(p \to q) = 1 \text{ et } \Pi(p) \geqslant \beta \text{ alors } \Pi(q) \geqslant \beta 
        \end{equation}
        \item \emph{Modus Tollens possibiliste} : 
        \begin{equation}
            \text{ Si } N(p \to q) \geqslant \alpha \text{ et } \Pi(q) \leqslant \beta \text{ alors } \Pi(p) \leqslant \max(1 - \alpha, \beta)  
        \end{equation}
    \end{itemize}
\end{proposition}

\begin{example}[Application du théorème fondamental]
    On veut savoir si l'on peut déduire que l'individu ne vole pas avec une certitude $\alpha = 0.9$. 
    On extrait la coupure $\Sigma_{0.9} = \{ M, M \to O, M \to \neg V \}$.
    Dans cette coupure, on a $(\Sigma_{0.9})^* \vdash \neg V$ par simple Modus Ponens classique.
    Ainsi, $N(\neg V) \geqslant 0.9$ est une déduction valide.
\end{example}

\subsection{Inconsistance}

En logique classique, si une base de connaissances admet une contradiction, alors on ne peut plus 
l'utiliser. En logique possibiliste, nous allons attribuer un \emph{niveau de confiance} à nos 
bases possibilistes appelé \emph{degré d'inconsistance}. 

\begin{definition}[Degré d'inconsistance]
    Soit $ \Sigma$ une base possibiliste. On définit son \emph{degré d'inconsistance}, noté 
    $Inc( \Sigma)$, par : 
        \begin{equation}
            \boxed{Inc( \Sigma) := \max \{ \alpha \mid \Sigma_\alpha \vdash \bot\}}
        \end{equation}
    Dans le cas où $ Inc( \Sigma) = 0$, on dit que la base est consistante. Si $Inc( \Sigma) = 1$, alors 
    elle est totalement inconsistante, on ne peut pas travailler avec. 
    Enfin, si $ 0 < Inc( \Sigma) < 1$, alors il y a un conflit entre des informations incertaines. 
\end{definition}

\begin{proposition}
    Ce niveau de confiance en une base possibiliste permet d'orienter le raisonnement définit plus 
    haut. En effet, pour qu'une conclusion $ \varphi$ soit considérée comme "valide", il sera 
    nécessaire que son degré de nécessité soit supérieur au niveau d'inconsistance de la base. 
        $$ \text{i.e} \quad N( \varphi) > Inc( \Sigma) $$  
\end{proposition}


\begin{example}[Calcul du degré d'inconsistance]
    Dans notre base $\Sigma$ augmentée du fait que l'individu est un manchot ($M, 1$):
    \begin{itemize}
        \item Une chaîne mène à $\neg V$ avec un degré $0.9$.
        \item Une autre mène à $V$ via $O$ avec un degré $\min(1, 1, 0.8) = 0.8$.
    \end{itemize}
    Le conflit entre $V$ et $\neg V$ au niveau le plus haut est $0.8$.
    On a donc $Inc(\Sigma) = 0.8$. 
    La conclusion $V$ ($0.8$) est rejetée, car $0.8 \ngtr 0.8$, tandis que $\neg V$ ($0.9$) est acceptée car $0.9 > 0.8$.
\end{example}

\subsection{Calcul par Résolution et Réfutation}

Pour automatiser le raisonnement sans tester tous les mondes possibles, on utilise la forme clausale 
(CNF) où chaque formule est une disjonction de littéraux.

\begin{definition}[Règle de Résolution Possibiliste (RP)] 
    Soient deux clauses possibilistes $(\lnot p \lor q, \alpha)$ et $(p \lor r, \beta)$. 
    La règle de résolution permet d'engendrer un \emph{résolvant} : 
    \begin{equation} 
        (\neg p \vee q, \alpha), (p \vee r, \beta) \vdash_{RP} (q \vee r, \min(\alpha, \beta)) 
    \end{equation} 
\end{definition}

Comme en logique classique, la résolution n'est pas complète pour la déduction directe, on utilise donc la \textbf{réfutation}.

\begin{theorem}[Complétude de la résolution] 
    $ \Sigma \vdash_\alpha \varphi$  si et seulement si on peut déduire la clause vide $ \bot$  avec un 
    poids $ \alpha$ à partir de la base augmentée de la négation de la conclusion : 
    \begin{equation} 
        \Sigma \cup { (\neg \varphi, 1) } \vdash_{RP} (\perp, \alpha) 
    \end{equation} 
    La valeur $ \alpha$ obtenue correspond au degré d'inconsistance de la base augmentée.
\end{theorem}

\begin{example}[Manchot et Vol] 
    Reprenons la base $ \Sigma = \{(M, 1), (\lnot M \lor O, 1), (\lnot M \lor \lnot V, 0.9), (\lnot O \lor V, 0.8)\}$
    \begin{itemize} 
        \item Le degré d'inconsistance de la base seule est $Inc( \Sigma) = 0.8$ (conflit sur V). 
        \item Pour prouver $\lnot V$, on ajoute $(V,1)$. On déduit $(\bot,0.9)$ par résolution. 
        \item Comme $0.9 > 0.8$, la conclusion "l'individu ne vole pas" est validée avec $ N(\lnot V) \geqslant 0.9$. 
    \end{itemize}
\end{example}






